\chapter{Discrete setting on AMR}\label{ch:discrete-setting}
Thesis Chapters~2, 6, and~7 introduce control-volume sets, function spaces,
and compatible operators~\cite{remmerswaal2023}. This chapter will extend
that setting to hanging faces, dual volumes, and vertex patches on the
block tree.

\section{Control volumes and functions}
The sets $\cvset{C}$ and $\cvset{F}$ denote leaf cells and their faces;
$\funspace{C}$ and $\funspace{F}$ are the respective function spaces.
\begin{implnote}
\apitype{CellField}, \apitype{FaceField}, \apitype{FaceVelocity}.
\end{implnote}

\section{Hanging faces and dual volumes}
The incidence fraction $\alphacf$ describes the part of cell $c$ associated
with face $f$; $\omega_f$ denotes the staggered control volume.
\begin{implnote}
\apitype{BlockTreeMesh}, \apitype{StaggeredVelocityAdvectionWorkspace}.
\end{implnote}

\section{Compatible operators}
The discrete divergence $\divh$ and gradient $\gradh$ and their AMR
adjointness and interpolation factorisation will be specified here.
\begin{implnote}
\apitype{GradientOperator}, \apifun{divergence!}, \apifun{gradient!}.
\end{implnote}
